\section{Diskrete Formulierung der Alignment-Optimierung}
\subsection{Diskretisierung des Gebiets}
\begin{definition}[Diskretisierung]
Es sei
\[
0=s_0<s_1<\dots<s_N=L
\]
eine Zerlegung von \([0,L]\).
\end{definition}
Wir bezeichnen \(\Delta s_i=s_{i+1}-s_i\).

\subsection{Diskrete Variablen}
\begin{definition}[Diskrete Krümmung]
\[\kappa_i\approx\kappa(s_i)\]
bezeichnet die Krümmungswerte an den Knoten.
\end{definition}
\begin{definition}[Diskrete Überhöhung]
\[\phi_i\approx\phi(s_i)\]
bezeichnet die Überhöhungswerte.
\end{definition}
\begin{definition}[Parametervektor]
Die Optimierungsvariable ist
\[
\mathbf x=(\kappa_0,\dots,\kappa_N,\phi_0,\dots,\phi_N)\in\R^{2(N+1)}.
\]
\end{definition}

\subsection{Diskrete Pose-Rekonstruktion}
Der pose3-Zustand wird numerisch integriert.
\begin{definition}[Diskrete Integration]
Aus den Anfangsbedingungen \(\gamma_0,\theta_0,\phi_0\) berechnen wir iterativ
\[
\theta_{i+1}=\theta_i+\kappa_i\Delta s_i,
\]
\[
\gamma_{i+1}=\gamma_i+\Delta s_i(\cos\theta_i,\sin\theta_i),
\]
\[
\phi_{i+1}=\phi_i+\omega_i\Delta s_i.
\]
\end{definition}
Dabei kann \(\omega_i\) durch Differenzen von \(\phi_i\) ausgedrückt werden.
Diese Rekonstruktion entspricht üblichen numerischen Integrationsideen
\cite{Hairer1993ODE,Press2007NumericalRecipes}.

\subsection{Finite Differenzen}
\begin{definition}[Erste Ableitung]
\[
\kappa'(s_i)\approx\frac{\kappa_{i+1}-\kappa_i}{\Delta s_i},
\quad
\phi'(s_i)\approx\frac{\phi_{i+1}-\phi_i}{\Delta s_i}.
\]
\end{definition}
\begin{definition}[Zweite Ableitung]
\[
\kappa''(s_i)\approx
\frac{\kappa_{i+1}-2\kappa_i+\kappa_{i-1}}{(\Delta s_i)^2}.
\]
\end{definition}

\subsection{Diskrete Zielfunktionen}
\subsubsection{Krümmungsglattheit}
\[
J_\kappa(\mathbf x)=\sum_{i=0}^{N-1}
\left(\frac{\kappa_{i+1}-\kappa_i}{\Delta s_i}\right)^2\Delta s_i.
\]
\subsubsection{Überhöhungsglattheit}
\[
J_\phi(\mathbf x)=\sum_{i=0}^{N-1}
\left(\frac{\phi_{i+1}-\phi_i}{\Delta s_i}\right)^2\Delta s_i.
\]
\subsubsection{Dynamische Zielfunktion}
\[
J_{\mathrm{dyn}}(\mathbf x)=\sum_{i=0}^{N}
\bigl(v^2\kappa_i-g\sin\phi_i\bigr)^2\Delta s_i.
\]
\subsubsection{Schwerpunkt-Zielfunktion}
Mit finiten Differenzen approximieren wir \(\|\gamma_c''(s_i)\|^2\) durch
Differenzen zweiter Ordnung:
\[
J_{\mathrm{veh}}(\mathbf x)=\sum_{i=1}^{N-1}
\|\gamma_{c,i+1}-2\gamma_{c,i}+\gamma_{c,i-1}\|^2.
\]

\subsection{Kombiniertes diskretes Problem}
\begin{definition}[Diskrete Zielfunktion]
\[
J(\mathbf x)=\alpha J_{\mathrm{veh}}+\beta J_{\mathrm{dyn}}
+\gamma J_\kappa+\delta J_\phi.
\]
\end{definition}

\subsection{Bedingungen}
\begin{definition}[Gleichheitsbedingungen]
\[\kappa_0=0,\quad\kappa_N=\kappa_1.\]
\end{definition}
\begin{definition}[Ungleichheitsbedingungen]
\[|\kappa_i|\leq\kappa_{\max},\quad|\phi_i|\leq\phi_{\max}.\]
\end{definition}
Weitere Bedingungen können Ableitungsgrenzen, Stetigkeitsbedingungen und
Alignment-Ankerpunkte umfassen.

\subsection{Abschließendes Optimierungsproblem}
\begin{definition}[Diskretes Optimierungsproblem]
Gesucht ist \(\mathbf x^\ast\in\R^{2(N+1)}\) mit
\[
J(\mathbf x^\ast)=\min_{\mathbf x}J(\mathbf x)
\]
unter den gegebenen Bedingungen.
\end{definition}

\subsection{Struktur des Problems}
Das Problem besitzt die Struktur eines nichtlinearen
Kleinste-Quadrate-Problems und eignet sich für SQP, Gauss--Newton-Verfahren und
dünn besetzte Optimierungstechniken.

\subsection{Verbindung zur Implementierung}
Die diskrete Formulierung entspricht unmittelbar einem Rechenmodell:
\(\kappa_i\) entsprechen Elementparametern, \(\phi_i\) Überhöhungsstichproben,
und die Rekonstruktion entspricht der numerischen Softwarebewertung.

\begin{summary}
Das kontinuierliche Alignment-Optimierungsproblem lässt sich auf ein
endlichdimensionales nichtlineares Optimierungsproblem reduzieren. Dies
verbindet theoretische Modellierung unmittelbar mit praktischer Implementierung.
\end{summary}
